\section{谱图景}
\markboth{谱图景}{}

\ob{4.1}{blk:4.1}\textbf{谱 Eisenstein 级数。}考虑叠的映射
\[
  \LocSys_{\Ch G}\xleftarrow{\fP\st}\LocSys_{\Ch B}\xrightarrow{\fQ\st}\LocSys_{\Ch T}.
\]
谱 Eisenstein 级数函子
\[
  \Eis\st : \QCoh(\LocSys_{\Ch T})\to \QCoh_N(\LocSys_{\Ch G})
\]
定义为复合 $(\fP\st)_*\circ(\fQ\st)^*$。这里 $\QCoh_N(\LocSys_{\Ch G})$ 是 $\QCoh(\LocSys_{\Ch G})$ 的放大，由 [GL:formulation] 引入。

把 $\Eis\st$ 的右伴随记作 $\CT\st$，称之为谱常数项函子。\footnote{为使 $\CT\st$ 作为保余极限的函子存在，需要 $\Eis\st$ 把紧对象送到紧对象。这正是引入修改 $\QCoh(\LocSys_{\Ch G})\to\QCoh_N(\LocSys_{\Ch G})$ 的原因：我们需要把 $\QCoh(\LocSys_{\Ch G})$ 的紧对象从「只是完美复形」扩大到包含 $\Eis\st$ 的像中的那些凝聚复形。}

几何朗兰兹等价所应具备的性质之一是：它把函子 $\Eis_!$ 与（按 $\rho(\omega_X)$ 平移后的）$\Eis\st$ 相缠，也就是说，下图必须交换：
\begin{equation}
\begin{tikzcd}[row sep=3.4em, column sep=3.6em, ampersand replacement=\&]
  D(\Bun_T) \& \QCoh(\LocSys_{\Ch T})
    \arrow[l, "-\rho(\omega_X)\circ\Psi_T"'] \\
  D(\Bun_G) \& \QCoh_N(\LocSys_{\Ch G})
    \arrow[l, "\Psi_G"']
    \arrow[from=1-1, to=2-1, "\Eis_!"]
    \arrow[from=1-2, to=2-2, "\Eis\st"']
\end{tikzcd}
\label{eq:4.1}
\end{equation}
这里 $\Psi_G$ 记 $G$ 的猜想性几何朗兰兹函子，而 $\Psi_T$ 是实现 $T$ 的几何朗兰兹对应的 Fourier-Mukai 变换。

本节的目标——本节遵循 [GL:Eis]——是解释范畴 $D(\Bun\rt_T)$ 与函子 $\Eis\rt$ 的谱意义。

\ob{4.2}{blk:4.2}\textbf{带纵向联络的拟凝聚层。}再次考虑映射
\[
  \fP\st : \LocSys_{\Ch B}\to\LocSys_{\Ch G},
\]
并令 $T_{\LocSys_{\Ch B}/\LocSys_{\Ch G}}$ 为 $\LocSys_{\Ch B}$ 上关于 $\fP\st$ 的纵向量场代数胚。

考虑范畴
\[
  T_{\LocSys_{\Ch B}/\LocSys_{\Ch G}}\text{-}\mathrm{mod}^r,
\]
即 $\QCoh(\LocSys_{\Ch B})$ 中的右 $T_{\LocSys_{\Ch B}/\LocSys_{\Ch G}}$-模构成的范畴。

泛包络代数 $U\bigl(T_{\LocSys_{\Ch B}/\LocSys_{\Ch G}}\bigr)$ 可以看作 $\QCoh(\LocSys_{\Ch B})$ 上的一个单子
\[
  \QCoh(\LocSys_{\Ch B})\to\QCoh(\LocSys_{\Ch B}),
\]
它等于遗忘函子
\[
  T_{\LocSys_{\Ch B}/\LocSys_{\Ch G}}\text{-}\mathrm{mod}^r\to\QCoh(\LocSys_{\Ch B})
\]
与其左伴随（即自由函子）的复合。

\begin{annote}
原文此句末作「i.e., the forgetful functor」，与前半句自相矛盾：遗忘函子的左伴随是自由函子 $\mathcal F\mapsto U(T)\otimes\mathcal F$，单子 $U(T)\otimes-$ 是「遗忘 $\circ$ 自由」，而不是遗忘函子本身。译文按「自由函子」译出，原文照录于工程文件疑点清单。\par\vspace{0.15em}{\raggedleft 【来源：原书 p.6 §4.2（A）；译者按（单子与伴随函子的一般理论）】\par}
\end{annote}

\ob{4.3}{blk:4.3}\textbf{精细谱 Eisenstein 级数。}函子
\[
  \fP\st_* : \QCoh(\LocSys_{\Ch B})\to\QCoh_N(\LocSys_{\Ch G})
\]
典范地分解为
\[
  \QCoh(\LocSys_{\Ch B})\xrightarrow{U(T_{\LocSys_{\Ch B}/\LocSys_{\Ch G}})}
  T_{\LocSys_{\Ch B}/\LocSys_{\Ch G}}\text{-}\mathrm{mod}^r
  \xrightarrow{\fP\st_{*,DR}}
  \QCoh_N(\LocSys_{\Ch G}).
\]
由于稍后将要解释的理由，我们把上述函子 $\fP\st_{*,DR}$ 称作「精细谱 Eisenstein 级数」，并记作 $\Eis^{\st,ref}$。

\ob{4.4}{blk:4.4}\textbf{精细谱常数项。}现在考虑 $\fP\st_*$ 的右伴随，我们把它记作 $(\fP\st_N)^!$。注意 $(\fP\st_N)^!$ 是如下复合：
\[
  \QCoh_N(\LocSys_{\Ch G})\leftarrow \IndCoh(\LocSys_{\Ch G})
  \xleftarrow{(\fP\st_N)^!}\IndCoh(\LocSys_{\Ch B})
  \to \QCoh(\LocSys_{\Ch B}),
\]
其中最后一个箭头是包含函子
\[
  \QCoh(\LocSys_{\Ch B})\hookrightarrow\IndCoh(\LocSys_{\Ch B})
\]
的右伴随。

函子 $(\fP\st_N)^!$ 典范地分解为
\[
  \QCoh_N(\LocSys_{\Ch G})\to
  T_{\LocSys_{\Ch B}/\LocSys_{\Ch G}}\text{-}\mathrm{mod}^r\to
  \QCoh(\LocSys_{\Ch B}),
\]
其中第二个箭头是遗忘函子。

我们把所得的函子 $\QCoh_N(\LocSys_{\Ch G})\to T_{\LocSys_{\Ch B}/\LocSys_{\Ch G}}\text{-}\mathrm{mod}^r$ 记作 $\CT^{\st,ref}$，并称之为「精细谱常数项」函子。函子对 $(\Eis^{\st,ref},\CT^{\st,ref})$ 自然地构成伴随对。

\ob{4.5}{blk:4.5}\textbf{谱侧的 Eisenstein 级数之常数项。}现在考虑单子
\[
  (\fP\st_N)^!\circ\fP\st_* : \QCoh(\LocSys_{\Ch B})\to\QCoh(\LocSys_{\Ch B}),
\]
它控制
\[
  \Hom_{\QCoh_N(\LocSys_{\Ch G})}\bigl(\fP\st_*(\mathcal F_1),\fP\st_*(\mathcal F_2)\bigr),
  \qquad \mathcal F_1,\mathcal F_2\in\QCoh(\LocSys_{\Ch B}).
\]

几乎是同义反复地，我们有：

\begin{othm}{引理}{4.5.1}{thm:4.5.1}
函子 $(\fP\st_N)^!\circ\fP\st_*$ 带有按元素 $w\in W$ 编号的典范滤过，使得：
\begin{enumerate}[label=(\alph*), leftmargin=2.4em, itemsep=0.25em]
\item 首项 $\bigl((\fP\st_N)^!\circ\fP\st_*\bigr)_{unit}$ 带有单子结构，且作为单子典范同构于 $U\bigl(T_{\LocSys_{\Ch B}/\LocSys_{\Ch G}}\bigr)$；
\item 末项 $\bigl((\fP\st_N)^!\circ\fP\st_*\bigr)_{w_0}$ 典范同构于复合
\[
  \QCoh(\LocSys_{\Ch B})\xrightarrow{\iota^*}\QCoh(\LocSys_{\Ch T})
  \xrightarrow{w_0}\QCoh(\LocSys_{\Ch T})
  \xrightarrow{\iota_*}\QCoh(\LocSys_{\Ch B}),
\]
其中 $\iota$ 是自然映射 $\QCoh(\LocSys_{\Ch T})\to\QCoh(\LocSys_{\Ch B})$。
\end{enumerate}
\end{othm}

\begin{annote}
引理 4.5.1 与命题 3.3.1 是同一现象的谱侧与几何侧：两侧的滤过都按 $w\in W$ 编号，首项都是「单子」（分别是 $U(T)$ 与 $\widetilde\Omega$），末项都是 $w_0$ 的作用。§4.6 的拟定理正是断言这两套滤过在 $\Psi_T$ 下逐项相合。\par\vspace{0.15em}{\raggedleft 【来源：原书 p.5 命题 3.3.1 与 p.7 引理 4.5.1 的对照（A）】\par}
\end{annote}

\ob{4.6}{blk:4.6}\textbf{几何图景与谱图景的关系。}回忆 $\fQ\st$ 记投影 $\LocSys_{\Ch B}\to\LocSys_{\Ch T}$。令 $\Psi_T$ 记 Fourier-Mukai 变换
\[
  D(\Bun_T)\simeq\QCoh(\LocSys_{\Ch T}).
\]

[GL:Eis] 的主旨如下：

\begin{othm}{拟定理}{4.6.1}{thm:4.6.1}
存在范畴的典范等价
\[
  \Psi_B : T_{\LocSys_{\Ch B}/\LocSys_{\Ch G}}\text{-}\mathrm{mod}^r
  \xrightarrow{\ \sim\ } D(\Bun\rt_T),
\]
它使下图交换：
\[
\begin{tikzcd}[row sep=3.6em, column sep=3.2em, ampersand replacement=\&]
  D(\Bun\rt_T) \&[-2em] T_{\LocSys_{\Ch B}/\LocSys_{\Ch G}}\text{-}\mathrm{mod}^r
    \arrow[l, "\Psi_B"'] \\
  D(\Bun_T) \& \QCoh(\LocSys_{\Ch T})
    \arrow[l, "-\rho(\omega_X)\circ\Psi_T"']
    \arrow[from=1-1, to=2-1, "(\fJ\rt_T)_!"]
    \arrow[from=1-2, to=2-2, "U(T_{\LocSys_{\Ch B}/\LocSys_{\Ch G}})\circ(\fQ\st)^*"]
\end{tikzcd}
\]
其中 $\rho(\omega_X)$ 是按点 $\rho(\omega_X)\in\Bun_T$ 平移的函子。
\end{othm}

注意，定理 4.6.1 可以改述如下。考虑单子
\[
  (\fQ\st)_*\circ U\bigl(T_{\LocSys_{\Ch B}/\LocSys_{\Ch G}}\bigr)\circ(\fQ\st)^*
  : \QCoh(\LocSys_{\Ch T})\to\QCoh(\LocSys_{\Ch T}),
\]
以及
\[
  \rho(\omega_X)\circ(\fJ\rt_T)_!\circ(\fJ\rt_T)^!\circ-\rho(\omega_X)
  =\rho(\omega_X)\circ\widetilde\Omega\circ-\rho(\omega_X)
  : D(\Bun_T)\to D(\Bun_T).
\]
断言是：$\Psi_T$ 把这两个单子相缠。

\begin{annote}
上述两图与两段改述里的 $\rho(\omega_X)$ 都指「按点平移」的函子，而不是乘以线丛。它与 $\widetilde\Omega$ 的复合放在 $\Psi_T$ 两侧，是为了把 §2.2.2 中「在 $D(\Bun_T)$ 上、未平移」的单子搬到与谱侧 $U(T)$ 对齐的位置；这正是几何朗兰兹里「用 $\rho$ 作中心平移」的常见做法。\par\vspace{0.15em}{\raggedleft 【来源：原书 p.8 §4.6（A）；[GL:Eis] 的做法（A）】\par}
\end{annote}

\ob{4.7}{blk:4.7}\textbf{小结。}于是，若我们已有满足 (4.1) 的函子 $\Psi_G$，则下述图也必须交换：
\begin{equation}
\begin{tikzcd}[row sep=3.6em, column sep=3.6em, ampersand replacement=\&]
  D(\Bun\rt_T) \& T_{\LocSys_{\Ch B}/\LocSys_{\Ch G}}\text{-}\mathrm{mod}^r
    \arrow[l, "-\rho(\omega_X)\circ\Psi_T"'] \\
  D(\Bun_G) \& \QCoh_N(\LocSys_{\Ch G})
    \arrow[l, "\Psi_G"']
    \arrow[from=1-1, to=2-1, "\Eis\rt"]
    \arrow[from=1-2, to=2-2, "\Eis^{\st,ref}"']
\end{tikzcd}
\label{eq:4.2}
\end{equation}

\begin{annote}
(4.2) 相对 (4.1) 的进步在于：把 (4.1) 中「谱侧的 $\Eis\st$」换成了「先 $U(T)$、再 $\Eis^{\st,ref}$」的分解，于是左边也相应地把 $D(\Bun_T)$ 换成 $D(\Bun\rt_T)$、把 $\Eis_!$ 换成 $\Eis\rt$。这一步之所以可能，靠的正是拟定理 4.6.1：它把 $T\text{-}\mathrm{mod}^r$ 与 $D(\Bun\rt_T)$ 等同起来。\par\vspace{0.15em}{\raggedleft 【来源：原书 p.8 §4.7（A）】\par}
\end{annote}
